\documentclass[conference]{IEEEtran}
\usepackage{flushend}
\usepackage{xurl}
\setlength{\emergencystretch}{2em}
\title{Sparse Freeway State Estimation Through Detector Outages}
\author{Arjun Vijay Prakash \\ \small IEEE ITSC 2027 | Regular paper working draft}

\begin{document}

\maketitle
\begin{abstract}
Persistence had lower speed and flow mean absolute error (MAE) than graph Huber reconstruction
in both directions of archived PeMS transfer under 30- and 60-minute detector outages.
On NGSIM trajectory aggregates, graph Huber improved speed recovery from I-80 to US-101,
but trailed the selected baseline in the reverse direction. A separate sparse storage model
became feasible when it allowed bounded tracking discrepancies. The PeMS readiness check
qualified zero of the ten days required for a frozen evaluation.
\end{abstract}

\section*{Problem and estimator}
The observation grid keeps speed, flow, counts, occupancy, and storage in their original
units. It records missing values, arrival times, and topology versions. Queue forecasts use
only readings that had arrived by issuance.

The graph smoother uses a tree prior and penalizes differences across neighboring cells
and times. A separate sparse model sets each change in storage equal to the integrated
net flow, with distinct internal, inlet, outlet, and ramp flows. Its point fit uses soft
observation residuals. A linear program uses hard observation intervals to bound selected
quantities and reports infeasibility when those intervals conflict.

\section*{Archived transfer}
PeMS training and calibration used first-day blocks separated by 30 minutes. The other
corridor supplied the next day's evaluation windows. All methods scored the same targets.
NGSIM used complete destination recordings at one-minute cadence. US-101 training and
calibration blocks were six minutes apart within one recording.

\begin{table}[h]
\centering\small
\caption{Contiguous-mask speed MAE in km/h. Reference selected on calibration.}
\begin{tabular}{lrrrr}
\hline
Transfer & min & $n$ & Huber & Ref. \\
\hline
I-880 to I-10 & 30 & 24 & 12.13 & 3.56 \\
 & 60 & 48 & 14.14 & 4.21 \\
I-10 to I-880 & 30 & 30 & 22.41 & 2.20 \\
 & 60 & 60 & 24.35 & 2.07 \\
I-80 to US-101 & 4 & 12 & 6.11 & 10.28 \\
 & 8 & 24 & 7.53 & 10.43 \\
US-101 to I-80 & 4 & 24 & 7.41 & 6.95 \\
 & 8 & 48 & 8.86 & 6.38 \\
\hline
\end{tabular}
\end{table}

Calibration selected persistence for both PeMS directions and NGSIM US-101, and graph
recurrent for NGSIM I-80. Nominal 90\% PeMS flow intervals covered zero of 162 target
instances. Each PeMS direction supplied one paired group. Reverse NGSIM transfer had two.

\section*{Storage and boundary assumptions}
The strict US-101 storage chain was infeasible with track-derived boundary counts. A
discrepancy cap of 17 vehicles per cell and step, calibrated on I-80, restored feasibility.
Storage MAE was 6.25 vehicles versus 6.36 for persistence.

\section*{Matched development ablations}
Hard observation intervals failed in four biased-sensor cases. The scored fallback was
persistence. Tighter boundary limits cut mean inlet interval width from 2,454 to 326
vehicles and covered the simulated total in all eight feasible cases. The speed smoother
trailed the tree on contiguous masks under light and surge demand. Its nominal 90\%
intervals covered 66\% and 59\% of targets.

\section*{Planned frozen evaluation}
The frozen protocol needs ten eligible days on held-out stations. It scores 30- and
60-minute outages with paired day-level intervals and prespecified flow and random-mask
margins. A separate 4- and 8-minute trajectory protocol requires an independent archive.
Both runners bind code, input content, and configuration before evaluation.

\section*{References}
[1] Federal Highway Administration. Next Generation Simulation vehicle trajectories and
supporting data.
\url{https://catalog.data.gov/dataset/next-generation-simulation-ngsim-vehicle-trajectories-and-supporting-data}

[2] Caltrans. Performance Measurement System. \url{https://pems.dot.ca.gov/}

[3] Development evaluation and receipts.
\url{https://github.com/ArjunCodess/CREST/blob/main/docs/crest-evaluation.md}
\end{document}
